%%%%%%%%%%%%%%%%%%%%%%%%%%%%%%%%%%%%%%%%%%%%%%%%%%%%%%%%%%%%%%%%%%%%%%%%%%%%%%%%
\documentclass{DG-report}
%%%%%%%%%%%%%%%%%%%%%%%%%%%%%%%%%%%%%%%%%%%%%%%%%%%%%%%%%%%%%%%%%%%%%%%%%%%%%%%

\usepackage[czech]{babel}
\usepackage[utf8]{inputenc}
% \usepackage[utf8]{inputenc}
% \usepackage[cp1250]{inputenc}

\usepackage{xspace}

\lhead[\fancyplain{} {\scriptsize M. Češka: Matematické a inženýrské metody
       pro vývoj spolehlivých a bezpečných paralelních a distribuovaných
       počítačových systémů}]
    {\fancyplain{} {\scriptsize M. Češka: Matematické a inženýrské metody pro
       vývoj spolehlivých a bezpečných paralelních a distribuovaných
       počítačových systémů}}

\newcommand{\datum}[3]{#1\,#2\,#3\xspace}

\newcommand{\note}[1]{[[ \emph{#1} ]]}
\renewcommand{\note}[1]{}

%%%%%%%%%%%%%%%%%%%%%%%%%%%%%%%%%%%%%%%%%%%%%%%%%%%%%%%%%%%%%%%%%%%%%%%%%%%%%%%
\begin{document}
%%%%%%%%%%%%%%%%%%%%%%%%%%%%%%%%%%%%%%%%%%%%%%%%%%%%%%%%%%%%%%%%%%%%%%%%%%%%%%%

%%%%%%%%%%%%%%%%%%%%%%%%%%%%%%%%%%%%%%%%%%%%%%%%%%%%%%%%%%%%%%%%%%%%%%%%%%%%%%%
\subsection*{Podklady studenta projektu pro průběžnou zprávu}
%%%%%%%%%%%%%%%%%%%%%%%%%%%%%%%%%%%%%%%%%%%%%%%%%%%%%%%%%%%%%%%%%%%%%%%%%%%%%%%

\note{Níže uveďte jméno, univerzitu, ročník, školitele a předpokládaný/finální
název disertace. Ti, kdo již obhájili, použijí makro \texttt{absolvent} namísto
\texttt{student}.}

%===============================================================================
\student{Jiří Slabý}		% \student nebo \absolvent
  {prof. RNDr. Antonín Kučera, Ph.D.}		% školitel
  {FI MU}				% univerzita
  {3.}					% ročník
  {Automatic Bug-finding Techniques for Linux Kernel}		% (předpokládaný) název disertace
%===============================================================================

%-------------------------------------------------------------------------------
\kratkeinfo
%-------------------------------------------------------------------------------

\note{Krátká charakteristika studia v daném roce (2--3 věty, nejlépe s odkazy na
publikace). Pište ve 3. osobě, prosím dodržte formát -- níže je uveden
příklad:}\bigskip

\noindent Student vedl jeden kurz na FI MUNI, účastnil se vybraných
seminářů na FI MU (včetně 1~aktivního vystoupení), složil státní
doktorskou zkoušku a obhájil teze své disertační práce. Ve výzkumu se zaměřil
na oblast statické analýzy jádra operačního systému Linux a 
vytvořil prototyp \cite{Slaby10a}.

%-------------------------------------------------------------------------------
\publikace
%-------------------------------------------------------------------------------

\note{Níže uveďte publikační a realizační výstupy za daný rok (uveďte i přijaté
publikace, s poznámkou přijato, příští rok dodáte definitivní záznam). Prosím
neužívejte slepě výpis z různých informačních systémů -- pokuste se dodržet
formát běžně generovaný z BiBTeX záznamů (bbl) tak, jak je použit níže.
Publikace označte jako Voprsalek10a, Voprsalek10b atd., ať se to dá nějak
rozumně sestavit za všechny studenty. \textbf{Vykazuje-li více studentů tutéž
publikaci, použijte totéž označení!!!!!!} Prázdné kategorie prosím
smažte.}\bigskip

%!!!!!!!!!!!!!!!!!!!!!!!!!!!!!!!!!!!!!!!!!!!!!!!!!!!!!!!!!!!!!!!!!!!!!

\renewcommand{\refname}{\vspace*{-12mm}}

\noindent {\bfseries Vytvořený software, hardware:}

\begin{thebibliography}{10}

\bibitem{Slaby10a} J.~Obdržálek, M.~Trtík, J.~Slabý. \textsc{Stanse}. Plně
funkční nástroj pro statickou analýzu jádra OS Linux dostupný na
\texttt{stanse.fi.muni.cz}.

\vspace{1.5mm}\hspace*{-9.3mm} \noindent {\bfseries Publikace ve sbornících
mezinárodních konferencí:}

\bibitem{Slaby10b}
J.~Škrabálek, L.~Tokárová, J.~Slabý, T.~Pitner
\newblock {Integrated Approach in Management and Design of ModernWeb-Based Services}.
\newblock In {\em Proc. of ISD'10}. Springer, 2010.

\bibitem{Slaby10c}
J.~Škrabálek, T.~Ludík, J.~Slabý, T.~Pitner
\newblock {Web-based Service for Collaborative Organization of Academic Events -- Case Study of ''Takeplace''}.
\newblock In {\em Proc. of SYNASC'10}. Springer, 2010.

\end{thebibliography}

%-------------------------------------------------------------------------------
\ucast
%-------------------------------------------------------------------------------

\note{Uveďte semináře, kterých jste se pravidelně nebo alespoň příležitostně
účastnili, resp. na kterých jste vystupovali. Uveďte i případná vystoupení na
navštívených zahraničních pracovištích. Dodržte prosím níže použitý
formát.}\bigskip

\bigskip

\noindent Pravidelná účast na seminářích:\begin{itemize}

  \item Formální metody v teorii i praxi na FI MU.
  
\end{itemize}

\noindent Příležitostná účast na seminářích:\begin{itemize}

  \item Seminar on Concurrency na FI MU.
  
\end{itemize}

\noindent Přednesené seminářové příspěvky:\begin{itemize}

  \item 3. 12. 2010: Nástroj Havoc, Seminar on Concurrency na FI MU.
  
\end{itemize}

%%%%%%%%%%%%%%%%%%%%%%%%%%%%%%%%%%%%%%%%%%%%%%%%%%%%%%%%%%%%%%%%%%%%%%%%%%%%%%%%
\subsection*{Zprávy ze zahraničních cest studentů hrazených z~projektu:}
%%%%%%%%%%%%%%%%%%%%%%%%%%%%%%%%%%%%%%%%%%%%%%%%%%%%%%%%%%%%%%%%%%%%%%%%%%%%%%%%

\note{Uveďte krátké zprávy ze svých zahraničních cest: zvláště krátké mohou být
u prezentace článku na nějaké konferenci, o něco delší by měly být zprávy z
letních škol a pracovních návštěv. Dodržte prosím následující formát (bude se
sestavovat do jednoho seznamu za celý projekt).}\bigskip

\begin{itemize}

  \item J.~Slabý, EmbeddedWorld 2010, Norimberk, Německo, 2.--3.3.2010,
  prezentace nástroje \cite{Slaby10a}.

\end{itemize}

%%%%%%%%%%%%%%%%%%%%%%%%%%%%%%%%%%%%%%%%%%%%%%%%%%%%%%%%%%%%%%%%%%%%%%%%%%%%%%

\label{KONEC}

%%%%%%%%%%%%%%%%%%%%%%%%%%%%%%%%%%%%%%%%%%%%%%%%%%%%%%%%%%%%%%%%%%%%%%%%%%%%%%%
\end{document}
%%%%%%%%%%%%%%%%%%%%%%%%%%%%%%%%%%%%%%%%%%%%%%%%%%%%%%%%%%%%%%%%%%%%%%%%%%%%%%%
